\section{Conclusion} 
Retrieval-augmented generation (RAG) demands a paradigm shift from relevance to utility—where passages are valued not merely for topical alignment but for their usefulness to enable complete, accurate answers. 
Utility-based RAG faces a critical scalability bottleneck: the high cost of using LLMs for utility judgments limits practical deployment to $\sim$20 passages per query, significantly hindering performance over large-scale passages.   
Since ranking requires a fixed strategy to determine passage count, we propose distilling the utility-based selection capability from larger models into smaller, efficient models. 
Key contributions and findings:
(1) We propose a novel utility-based selection distillation method that trains efficient student models (e.g., Qwen3-1.7B) to jointly learn pseudo-answer generation and utility assessment from teacher LLMs (Qwen3-32B). This bypasses the need for direct LLM inference during deployment. 
(2) We empirically validate that utility-driven passage selection (unlike traditional relevance ranking) is essential for complex QA tasks (e.g., HotpotQA), where synthesizing information across multiple documents is paramount. 
(3) Our distilled model employs a front-to-back sliding window strategy to adaptively select minimal high-utility passage sets from large-scale candidates (top-100), dynamically skipping low-utility regions. 
Adaptive selection yields higher-quality answers with fewer passages, accelerating inference. This reduces computational costs by 70\% while improving answer quality. 
Additionally, we will release the Qwen3-32B relevance ranking and utility-based selection annotations for the 100k MS MARCO dataset, providing a high-quality dataset for future research in relevance ranking and utility-based selection. 
For future work, extending utility distillation to tasks beyond QA (e.g., summarization, fact verification) is promising in the future. 



% While modern dense retrievers and LLM-powered re-rankers excel at identifying topically relevant passages, we demonstrated that selecting evidence for RAG solely based on relevance is insufficient for optimal RAG performance. 
% Relevance ranking overlooks the critical dimension of information utility – the actual usefulness of passages in generating accurate and comprehensive answers – and lacks query-specific adaptability in determining the optimal number of passages to provide to the generator.

% To address these limitations, we introduced and validated utility-based selection as a superior alternative to ranking for RAG. By enabling an LLM to generate a pseudo-answer and explicitly select documents most useful for answering the query itself, our approach provides a more flexible and effective mechanism for enhancing the information provided to the generator. Crucially, we mitigated the computational barrier of online utility-based selection by successfully distilling the capability from a powerful 32B LLM (Qwen3-32B) into a highly efficient 1.7B model (\utility$_{1.7B}$). 
% Our key findings are: 
% 1)Experiments on NQ and HotpotQA, particularly on complex queries, confirmed that relevance-centric approaches fall short. Utility-based selection, focused on answer generation needs, demonstrably improves RAG outcomes. 
% 2)  Utility-based selection proved more effective than traditional ranking followed by fixed-threshold filtering in identifying the most beneficial documents for the generator, offering inherent adaptability to query-specific information requirements. 
% To foster further research, we publicly release the Qwen-32B relevance ranking annotations for the MS MARCO dataset. Our work establishes utility-based selection as a promising direction for enhancing RAG efficiency and effectiveness, while the released resources and efficient distilled models (\rank$_{1.7B}$, \utility$_{1.7B}$) provide practical tools for the community.